\documentclass[paper=screen,mode=present,style=zysimple]{powerdot}

\usepackage{fancyvrb,verbatim,listings}
\usepackage{proof}
\usepackage[all]{xypic,xy}

\usepackage{latexsym,stmaryrd,amsmath,amsfonts}

%% \newtheorem{definition}{Definition}
%% \newtheorem{theorem}{Theorem}
%% \newtheorem{lemma}{Lemma}
\newtheorem{prop}{Proposition}
%% \newtheorem{corollary}{Corollary}
%% \newenvironment{proof}{{\sc Proof . }}{{\hfill $\square$}}
%% \newcommand{\myqed}{\hfill $\square$}
%% \newcommand{\ped}{\vspace{-1.6em}}
%% \newcommand{\note}[1]{\noindent\fbox{\parbox{\textwidth}{\em #1}}}

\newcommand{\abind}[2]{\mathtt{bind} \ (\, {#1} \,)  \ \mathtt{in} \ {#2}}
\newcommand{\acase}{\mathtt{case}}
\newcommand{\adv}{\mathcal{A}}
\newcommand{\aspecto}[1]{\xrightarrow{\; {#1} \;}}
\newcommand{\asplinear}{\mathfrak{l}}
\newcommand{\aspmodal}{\mathfrak{m}}
\newcommand{\aspnormal}{\mathfrak{n}}
%
\newcommand{\betareduce}{\leadsto_\beta}
\newcommand{\betareduceto}{\twoheadrightarrow_\beta}
\newcommand{\bit}{\mathtt{B}}
\newcommand{\bitone}{\mathtt{B}_1}
\newcommand{\bitzero}{\mathtt{B}_0}
%
\newcommand{\caseto}{\triangleright}
\newcommand{\cbind}[3]{\mathtt{bind} \ {#1} = {#2} \ \mathtt{in} \ {#3}}
\newcommand{\ccase}[4]{\mathtt{case} \; {#1} \; \mathtt{of} \; \cnil \caseto {#2}; \;
                                          \bitzero x_0 \caseto {#3}; \; \bitone x_1 \caseto {#4}}
\newcommand{\cequiv}{\simeq}
\newcommand{\cif}[3]{\mathtt{if} \; {#1} \; \mathtt{then} \; {#2} \; \mathtt{else} \; {#3}}
\newcommand{\ctrue}{\mathtt{true}}
\newcommand{\cfalse}{\mathtt{false}}
\newcommand{\cfix}{\mathtt{fix}\,}
\newcommand{\cfixin}[2]{\mathtt{fix}\; {#1}\; \mathtt{in} \; {#2}}
\newcommand{\cfold}{\mathtt{fold}\,}
\newcommand{\chd}{\mathtt{hd}}
\newcommand{\cinj}{\mathtt{inj}}
\newcommand{\clift}[1]{\lfloor{#1}\rfloor}
\newcommand{\clet}[3]{\mathtt{let} \ {#1} = {#2} \ \mathtt{in} \ {#3}}
\newcommand{\cletst}[3]{\mathtt{liftlet} \ {#1} = {#2} \ \mathtt{in} \ {#3}}
\newcommand{\cletrec}[3]{\mathtt{letrec} \ {#1} = {#2} \ \mathtt{in} \ {#3}}
\newcommand{\cnil}{\mathtt{nil}}
\newcommand{\conc}{\! \bullet \!}
\newcommand{\concat}{\pmb{\mathit{conc}}}
\newcommand{\context}{\mathbb{C}}
\newcommand{\cop}{\;\mathtt{op}\;}
\newcommand{\cpair}[1]{\langle{#1}\rangle}
\newcommand{\cproj}{\mathtt{proj}}
\newcommand{\crand}{\mathtt{rand}}
\newcommand{\crec}{\mathtt{rec}}
\newcommand{\croll}{\mathtt{roll}\,}
\newcommand{\crule}[1]{\mbox{{\textsl{#1}}}}
\newcommand{\cslrtype}{\Vdash}
\newcommand{\csucc}{\mathtt{S}}
\newcommand{\ctl}{\mathtt{tl}}
\newcommand{\cut}{\pmb{\mathit{cut}}}
\newcommand{\cval}{\mathtt{return}}
\newcommand{\cwhile}[2]{\mathtt{while}\; {#1} \; \mathtt{do} \; {#2} \; \mathtt{od}}
\newcommand{\cvalof}{\pmb{valueof}}
\newcommand{\cunfold}{\mathtt{unfold}\,}
\newcommand{\cunit}{\mathtt{unit}}
\newcommand{\cunroll}{\mathtt{unroll}\,}
\newcommand{\cupdate}{\pmb{update}}
\newcommand{\czero}{\mathtt{O}}
%
\newcommand{\deduces}{\vdash}
\newcommand{\defeq}{\stackrel{\mathrm{def}}{=}}
\newcommand{\delay}[1]{\left\lfloor{#1}\right\rfloor}
\newcommand{\distr}{\mathcal{D}}
%
\newcommand{\eagerreduce}{\xrightarrow{\mathrm{eager}}}
\newcommand{\Eagerreduce}{\stackrel{\mathrm{eager}}{\Reduce}}
\newcommand{\eqtest}{\stackrel{\mbox{\tiny ?}}{=}}
\newcommand{\Eval}[1]{{\left\llbracket{#1}\right\rrbracket}}
\newcommand{\evalcontext}{\mathbb{E}}
\newcommand{\Example}{\emgreen{\pmb{\textsc{Example}\;}}}
\newcommand{\expansion}[1]{\tilde{{#1}}}
%
\newcommand{\ford}{\mathrm{ord}}
%
\newcommand{\getval}{\leftarrow}
%
\newcommand{\head}{\pmb{\mathit{hd}}}
\newcommand{\termhole}{[\;]}
\newcommand{\hoaretriple}[3]{\{{#1}\} \; {#2} \; \{{#3}\}}
%
\newcommand{\imply}{\Rightarrow}
\newcommand{\indist}{\simeq}
\newcommand{\iteration}{\pmb{\mathit{it}}}
%
\newcommand{\lazyreduce}{\xrightarrow{\mathrm{lazy}}}
\newcommand{\Lazyreduce}{\lazyreduce\!\!\!\!\!\to}
\newcommand{\ldot}{\, . \, }
\newcommand{\linearto}{\multimap}
\newcommand{\logrel}{(\mathcal{R}_\tau)_{\mbox{$\tau$ type}}}
\newcommand{\lr}[1]{\mathcal{R}_{#1}}
\newcommand{\lub}{\sqcup}
%
\newcommand{\metalambda}{\underline{\lambda}}
\newcommand{\modal}{\square}
%
\newcommand{\nat}{\mathbb{N}}
\newcommand{\nReduce}{\nrightarrow\!\!\!\!\to}
\newcommand{\nreduceto}{\nrightsquigarrow}
\newcommand{\nullenv}{\varnothing}
%
\newcommand{\orc}{\mathcal{O}}
\newcommand{\oslrtype}{\Vdash}
%
\newcommand{\porder}{\sqsubseteq}
\newcommand{\prefix}{\pmb{\mathit{pref}}}
\newcommand{\proba}{\mathbf{Pr}}
%
\newcommand{\randlen}{\pmb{\mathit{rl}}}
\newcommand{\randnum}{\pmb{\mathit{rn}}}
\newcommand{\randstr}{\pmb{\mathit{rs}}}
\newcommand{\reduce}{\leadsto}
\newcommand{\reduceto}{\rightsquigarrow\hspace*{-6px}\shortrightarrow}
\newcommand{\Reduce}{\to\!\!\!\to}
\newcommand{\Reduceto}{\Downarrow}
%
\newcommand{\set}[1]{\{{#1}\}}
\newcommand{\setbits}{\mathbb{B}}
\newcommand{\slrtype}{\vdash}
\newcommand{\splt}{\pmb{\mathit{split}}}
\newcommand{\suffix}{\pmb{\mathit{suff}}}
\newcommand{\sysname}{\pmb{\mathcal{C}}}
%
\newcommand{\tail}{\pmb{\mathit{tl}}}
\newcommand{\tbits}{\mathsf{Bits}}
\newcommand{\tbool}{\mathsf{Bool}}
\newcommand{\tnat}{\mathsf{Nat}}
\newcommand{\tcomp}{\mathsf{T}}
\newcommand{\tlift}[1]{\lfloor{#1}\rfloor}
\newcommand{\tlist}{\mathsf{List}}
\newcommand{\tself}{\mathsf{self}}
\newcommand{\tstore}{\mathsf{Store}}
\newcommand{\tstream}{\mathsf{Stream}}
\newcommand{\tloc}{\mathsf{Loc}}
\newcommand{\ttcase}{\mathtt{case}}
\newcommand{\ttif}{\mathtt{if}}
\newcommand{\ttthen}{\mathtt{then}}
\newcommand{\ttelse}{\mathtt{else}}
\newcommand{\ttin}{\mathtt{in}}
\newcommand{\ttlet}{\mathtt{let}}
\newcommand{\ttof}{\mathtt{of}}
\newcommand{\tunit}{\mathsf{Unit}}
\newcommand{\Type}{\mathcal{T}\!\mathit{ype}}
%
\newcommand{\vbool}{\pmb{\mathbb{B}}}
\newcommand{\vff}{\mathfrak{f}\!\mathfrak{f}}
\newcommand{\vnat}{\pmb{\mathrm{N}}}
\newcommand{\vnil}{\epsilon}
\newcommand{\vtt}{\mathfrak{t}\!\mathfrak{t}}
%
\newcommand{\weakpre}{\pmb{wp}}




\pdsetup{lf={Theory and Practice of Functional Programming}}

\title{{\small \textsc{Theory and Practice of Functional Programming}}
\\[12pt]
{\Large \color{TitleColor} Imperative Languages and Monads}
}
\author{\small \textbf{Dr. ZHANG Yu}
  \\[2pt] \small Institute of Software, Chinese Academy of Sciences 
  \\[5pt] \includegraphics[height=.04\slideheight]{iscas.eps} 
  \qquad \qquad 
  \includegraphics[height=.04\slideheight]{cas.eps}
}
\date{\scriptsize Fall term, 2012 \qquad GUCAS, Beijing}

\newcommand{\RED}{\mathsf{RED}}

\begin{document}

\maketitle 

\begin{slide}{A toy imperative language WHILE}
WHILE is a simple imperative language for developing analysis technique for imperative programs.
\begin{itemize}
\item Expressions
\begin{itemize}
\item Integer expressions
\vspace*{-0.5em}
\[
E_n ::= x \mid 0 \mid 1 \mid \ldots \mid E_n \cop E_n 
\vspace*{-1.6em}
\]
\item Boolean expressions
\vspace*{-0.5em}
\[
E_b ::= \ctrue \mid \cfalse \mid E_n \eqtest E_n 
\vspace*{-0.6em}
\]
\end{itemize}
\item Programs 
\vspace*{-0.5em}
\[
\begin{array}{lrll}
P & ::= & x := E_n  & \mbox{assignment} 
\\[0.3em] & \mid & 
P; P & \mbox{sequential}
\\[0.3em] & \mid & 
\cif{E_b}{P}{P} & \mbox{conditional} 
\\[0.3em] & \mid & 
\cwhile{E_b}{P} & \mbox{while-loop}
\end{array}
\]
\end{itemize}
\end{slide}

\begin{slide}[method=direct]{Stores and locations}
\begin{itemize}
\item The execution of a WHILE program depends on the \emblue{store} (or memory) which feeds 
  the variables with concrete values. \\[0.3em]
\item A variable is essentially a \emblue{location} in the store.
\vspace*{-0.5em}
\[
x := y + 3
\vspace*{-0.6em}
\]
\begin{itemize}
\item Get the value from {\em location \verb|y|} and put it in a register;
\item Add 3 to the register;
\item Store the context of register into {\em location \verb|x|}.
\end{itemize}
\item L-value and R-value
\begin{itemize}
\item L-value: the location of the variable \\[0.2em]
\item R-value: the content in the store location of the variable
\end{itemize}
\end{itemize}
\end{slide}

\begin{slide}{Evaluation of expressions}
The evaluation of WHILE expressions depends on the store. 
\begin{itemize}
\item We simply regard variables in WHILE programs as locations. \\
\item \emblue{$\pmb{valueof}$}: 
  $\pmb{valueof}(S, x)$ is the value of the variable $x$ in the store $S$. \\
\item Evaluation of expressions: $\reduce \, \subseteq (Store \times Exp) \times Value$:
\[
\begin{array}{l}
(S, x) \reduce \pmb{valueof}(S, x) 
, \quad (S, n) \reduce n \\[0.3em]
(S, E_1 \cop E_2) \reduce n, 
\mbox{ if } (S, E_1) \reduce n_1, (S, E_2) \reduce n_2 \mbox{ and } n = n_1 \cop n_2 \\[0.3em]
(S, \ctrue) \reduce \vtt 
, \quad (S, \cfalse) \reduce \vff \\[0.3em]
(S, E_1 \eqtest E_2) \reduce \ctrue, \mbox{ if } (S, E_1) \reduce n, (S, E_2) \reduce n \\[0.3em]
(S, E_1 \eqtest E_2) \reduce \ctrue, \mbox{ if } (S, E_1) \reduce n_1, (S, E_2) \reduce n_1 \mbox{ and } n_1 \neq n_2 
\end{array}
\]
% The store $S$ is much like the environment in the denotational semantics of PCF.
\end{itemize}
\end{slide}

\begin{slide}{Small-step semantics}
Small-step operational semantics: 
% It is a relation between programs (w.r.t. stores): \vspace*{-0.5em}
\[
\Rightarrow \, \subseteq (Store \times Program) \times (Store \times Program) .
\vspace*{-1em}
\]
\begin{itemize}
\item It defines the execution of programs, {\em step by step}. \\
\item Final state $(S, \bullet)$: $\bullet$ stands for a null program. It also indicates \emred{termination}. \\
\item \emblue{$\pmb{update}(S, x, v)$}: a new store obtained from $S$ with the location $x$ containing value $v$. \\
\item Assignment: 
\[
\begin{prooftree}
(S, E) \reduce v  
\quad \pmb{update}(S, x, v) = S'
\justifies
(S, x := E) \Rightarrow (S', \bullet)
\end{prooftree}
\vspace*{-1.5em}
\]
\end{itemize}
\end{slide}

\begin{slide}{Small-step semantics}
\begin{itemize}
\item Conditional: 
\[\small
\begin{array}{c}
\begin{prooftree}
(S, E) \reduce \ctrue
\justifies
(S, \cif{E}{P_1}{P_2}) \Rightarrow (S, P_1)
\end{prooftree}
\quad
\begin{prooftree}
(S, E) \reduce \cfalse
\justifies
(S, \cif{E}{P_1}{P_2}) \Rightarrow (S, P_2)
\end{prooftree}
\end{array}
\vspace*{-1.5em}
\]
\item Sequential:
\[\small
\begin{prooftree}
(S, P_1) \Rightarrow (S', \bullet)
\justifies
(S, P_1 ; P_2) \Rightarrow (S', P_2)
\end{prooftree}
\vspace*{-1.5em}
\]
\item WHILE-loop:
\[\small
\begin{array}{c}
\begin{prooftree}
(S, E) \reduce \cfalse
\justifies
(S, \cwhile{E}{P}) \Rightarrow (S, \bullet)
\end{prooftree}
\\[1.8em]
\begin{prooftree}
(S, E) \reduce \ctrue
\justifies
(S, \cwhile E P) \Rightarrow (S, P; \cwhile{E}{P})
\end{prooftree}
\end{array}
\vspace*{-0.5em}
\]
The evaluation of WHILE-loop is a loop \emblue{unfolding} process. 
\end{itemize}
\end{slide}

\begin{slide}{Big-step semantics}
Big-step operational semantics:
$$
\Reduceto \subseteq (Store \times Program) \times Store
\vspace*{-1em}
$$
\begin{itemize}
\item It defines the final store after program execution (w.r.t. an initial store). \\
\item Assignment:
\[
\begin{prooftree}
(S, E) \reduce v  
\quad \pmb{update}(S, x, v) = S'
\justifies
(S, x := E) \Reduceto S'
\end{prooftree}
\]
\item Conditional:
\[
\begin{prooftree}
(S, E) \reduce \ctrue
\quad (S, P_1) \Reduceto S'
\justifies
(S, \cif{E}{P_1}{P_2}) \Reduceto S'
\end{prooftree}
\qquad 
\begin{prooftree}
(S, E) \reduce \cfalse
\quad (S, P_2) \Reduceto S'
\justifies
(S, \cif{E}{P_1}{P_2}) \Reduceto S'
\end{prooftree}
\]
\end{itemize}
\end{slide}

\begin{slide}{Big-step semantics}
% The big-step semantics
\begin{itemize}
\item Sequential:
\[
\begin{prooftree}
(S, P_1) \Reduceto S_1
\quad (S_1, P_2) \Reduceto S_2
\justifies
(S, P_1 ; P_2) \Reduceto S_2
\end{prooftree}
\vspace*{-0.5em}
\]
\item WHILE-loop:
\[
\begin{prooftree}
(S, E) \reduce \cfalse
\justifies
(S, \cwhile{E}{P}) \Reduceto S
\end{prooftree}
\]
\[
\begin{prooftree}
(S, E) \reduce \ctrue
\quad (S, P) \Reduceto S'
\quad (S', \cwhile{E}{P}) \Reduceto S''
\justifies
(S, \cwhile E P) \Reduceto S''
\end{prooftree}
\vspace*{-1.5em}
\]
\item Termination: \\
\begin{itemize}
\item Big-step evaluation implicitly assumes termination: 
  if $(S, P) \Reduceto S'$, then $P$ necessarily terminates.
\item Small-step evaluation does not assume termination.
\end{itemize}
\end{itemize}
\end{slide}

\begin{slide}[method=direct]{Example}
\scriptsize
\begin{verbatim}
        x := 5 ;
        y := 1;
        while (x > 0) do 
          y := x * y;
          x := x - 1;
        od 
\end{verbatim}
\normalsize
\begin{itemize}
\item Small-step evaluation
\vspace*{-0.6em}
\[\scriptsize
\begin{array}{rll}
 & (\set{x : 0, y : 0}, x := 5; y := 1; \cwhile{x > 0}{\ldots})
\\ \Rightarrow & 
(\set{x : 5, y : 0}, y := 1; \cwhile{x > 0}{\ldots})
\\ \Rightarrow & 
(\set{x : 5, y : 1}, \cwhile{x > 0}{y := y*x; x := x - 1})
\\ \Rightarrow & 
(\set{x : 5, y : 1}, y := y*x; x := x - 1; \cwhile{x > 0}{\ldots})
\\ \Rightarrow & \ldots \ldots
\end{array}
\vspace*{-1em}
\]
\item Big-step evaluation
\[\scriptsize
\hspace*{-2em}
\begin{prooftree}
\begin{prooftree}
\pmb{update}(\set{x:0, y:0}, x, 5) = \set{x: 5, y : 0}
\justifies
(\set{x : 0, y : 0}, x := 5) \Reduceto \set{x : 5, y : 0}
\end{prooftree}
\;
\begin{prooftree}
\cdots \quad \cdots
\justifies
(\set{x : 5, y : 0}, y := 1; \cwhile{\ldots}{\ldots}) \Reduceto \set{x : 0, y : 120}
\end{prooftree}
\justifies
(\set{x : 0, y : 0}, x := 5; y := 1; \cwhile{x > 0}{\ldots}) \Reduceto \set{x : 0; y : 120}
\end{prooftree}
\]
\normalsize
\end{itemize}
\end{slide}

\section{Encoding WHILE in PCF}

\begin{slide}{PCF with stores}
\begin{itemize}
\item Imperative languages can be encoded in purely functional languages. \\
\begin{itemize}
\item A better understanding of the relation between the two styles of languages. \\
\item This also helps to define denotational semantics for imperative languages.
\vspace*{-0.5em}
\end{itemize}
\item To interpret WHILE, we use an extended version of PCF \emblue{with stores}.
\[
\tau ::= \tnat \mid \tbool \mid \emblue{\tstore} \mid \emblue{\tloc} \mid \emred{\tlift{\tstore}}
\mid \tau \times \tau \mid \tau \to \tau
\]
$\tloc$ is the type for locations in a store (L-values of variables). \\[0.3em]
\end{itemize}
\end{slide}

\begin{slide}{PCF with stores}
\begin{itemize}
\item $\tlift{\tstore}$ is a \emblue{lifted type}, which contains all elements of $\tstore$ 
plus \emblue{non-termination}.
\vspace*{-0.3em}
\[
\begin{prooftree}
\Gamma \vdash e : \tstore 
\justifies
\Gamma \vdash \clift e : \tlift{\tstore}
\end{prooftree}
\qquad
\begin{prooftree}
\Gamma \vdash e : \tlift{\tstore}
\quad \Gamma, x: \tstore \vdash e' : \tlift\tstore  
\justifies
\Gamma \vdash \cletst{x}{e}{e'} : \tlift\tstore
\end{prooftree}
\vspace*{-1em}
\]
\begin{itemize}
\item The $\mathtt{letstore}$ expression: if $e$ does not terminate, 
  then the entire expression does not terminate either; 
  otherwise evaluates $e'$ with $x$ being the value of $e$. \\
\item $\tlift{\tstore}$ can be defined using sum and unit:  $\tlift{\tstore} \defeq \emblue{\tstore + \tunit}$.
\vspace*{-1em}
\end{itemize}
\item Pre-defined functions \vspace*{-0.5em}
\begin{itemize}
\item $\cvalof : \emblue{\tstore \times \tloc \to \tnat}$.\\
\item $ \cupdate : \emblue{\tstore \times \tloc \times \tnat \to \tstore}$. 
\vspace*{-1.6em}
\end{itemize}
\begin{eqnarray*}
\cvalof(\cupdate(s, l, v), l) & = & v
\\
\cupdate(s, l, \cvalof(s, l)) & = & s
\\
\cupdate(\cupdate(s, l, v_1), l, v_2) & = & \cupdate(s, l, v_2)
\\[-4em]
\end{eqnarray*}
\end{itemize}
\end{slide}

\begin{slide}{Interpretation}
\begin{itemize}
\item Interpretation of expressions $\Eval{\_}_E$: 
  PCF functions of \emblue{$\tstore \to \tnat$} (or \emblue{$\tstore \to \tbool$}).
\begin{eqnarray*}
\Eval{x} & \defeq & \lambda s : \tstore \ldot \cvalof(s, x) \\[0.3em]
\Eval{E_1 + E_2} & \defeq & \lambda s : \tstore \ldot \Eval{E_1} (s) + \Eval{E_2} (s) \\[0.3em]
\Eval{E_1 \eqtest E_2} & \defeq & 
\lambda s : \tstore \ldot \\
& & \qquad \cif{\Eval{E_1} (s) \eqtest \Eval{E_2} (s)}{\ctrue}{\cfalse} \\[0.3em]
\ldots & \ldots & 
\end{eqnarray*}
\end{itemize}
\end{slide}

\begin{slide}{Interpretation}
\begin{itemize}
\item Interpretation of programs $\Eval{\_}_P$: PCF functions of \emred{$\tstore \to \tlift\tstore$}.
\begin{itemize}
\item Assignment:
\vspace*{-0.5em}
\[
\Eval{x := E} \defeq \lambda s : \tstore \ldot \clift{\cupdate(s, x, \Eval{E} (s))}
\vspace*{-1.5em}
\]
\item Conditional: 
\vspace*{-0.5em}
\[
\begin{array}{l}
\Eval{\cif{E}{P_1}{P_2}} \defeq \\[0.2em]
\qquad\qquad \lambda s : \tstore \ldot \cif{\Eval{E}(s)}{\Eval{P_1} (s)}{\Eval{P_2} (s)}
\end{array}
\vspace*{-1em}
\]
\item Sequential: 
\vspace*{-0.5em}
\[
\Eval{P_1 ; P_2} \defeq \lambda s : \tstore \ldot 
\cletst{s_1}{\Eval{P_1} (s)}{\Eval{P_2} (s_1)}
\vspace*{-1.5em}
\]
\end{itemize}
\end{itemize}
\end{slide}

\begin{slide}{Interpretation of while-loop}
\begin{itemize}
\item The interpretation of while-loop is recursively defined.
\[
\begin{array}{l}
\Eval{\cwhile E P} =  
\\[0.2em] \qquad \qquad 
\lambda s : \tstore \ldot \ttif \; \Eval{E}(s) 
\\[0.2em] \hspace*{8.7em}
\ttthen \; \cletst{s'}{\Eval{P} (s)}{\Eval{\cwhile E P} (s')}
\\[0.2em] \hspace*{8.7em}
\ttelse \; \clift s 
\end{array}
\]
The definition with PCF fix-point operator:
\[
\begin{array}{l}
\Eval{\cwhile E P} =  
\\[0.2em] \qquad \, 
\cfix \, \lambda f : (\tstore \to \tlift\tstore) \ldot \lambda s : \tstore \ldot 
\\[0.2em] \qquad \qquad
\ttif \; {\Eval{E}(s)} \; 
\ttthen \; \cletst{s'}{\Eval{P} (s)}{f (s')} \; \ttelse \; \clift{s}
\end{array}
\]
\end{itemize}
\end{slide}

\begin{slide}{Example}
We can use PCF-store to reason about WHILE programs.
\begin{itemize}
\item Define $\pmb{skip} \defeq x := x$.
\vspace*{-0.7em}
\begin{eqnarray*}
\Eval{\pmb{skip}} & = & 
\lambda s : \tstore  \ldot \clift{\cupdate(s, x, \Eval x (s))} 
\\ & = & 
\lambda s : \tstore \ldot \clift{\cupdate(s, x, \cvalof(s, x))}
\\ & = & 
\lambda s : \tstore \ldot \clift{s}
\end{eqnarray*}
\vspace*{-2em}
\item Prove that $\Eval{\cwhile E P} = \Eval{\cif{E}{P; \cwhile E P}{\pmb{skip}}}$.
\vspace*{-0.7em}
\[
\begin{array}{ll}
& \Eval{\cwhile E P} 
\\[0.2em] = &
\cfix \, \lambda f \ldot \lambda s \ldot 
\cif{\Eval E (s)}{\cletst{s'}{\Eval P (s)}{f (s')}}{\clift s}
\\[0.2em] = & 
\lambda s \ldot \ttif \; \Eval E (s) \; \ttthen \; \cletst{s'}{\Eval P (s)}{\Eval{\cwhile E P}(s')}
\\ & \hspace*{1.8em}
\ttelse \; \clift s
\\[0.2em] = & 
\lambda s \ldot \cif{\Eval E (s)}{\Eval{P; \cwhile E P}(s)}{\Eval{\pmb{skip}(s)}}
\\[0.2em] = &
\Eval{\cif{E}{P; \cwhile E P}{\pmb{skip}}}
\end{array}
\]
\end{itemize}
\end{slide}

\begin{slide}{Equivalence between op. and de. semantics}
\begin{itemize}
\item For every WHILE-expression $E$, \emblue{$(S, E) \reduce v$ iff $\Eval E (S) = v$}.
\\[0.2em]
{\em Proof.} Done by induction on the expression structure.
\\[0.3em]
\item For every WHILE-program  $P$, \emblue{$(S, P) \Reduceto S'$ 
  iff $\Eval P (S) = \clift{S'}$ and $\Eval P (S) \neq \bot$}.
\\[0.2em]
{\em Proof.} This is also a proof by induction on program structure. But for the WHILE-loop, 
we need an extra induction on the unfolding length (for every terminating loop):
\[
\begin{array}{l}
\pmb{length} (S, \cwhile E P) \defeq n \mbox{ such that} 
\\[0.2em] \qquad \qquad 
S = S_0 
\\[0.2em] \qquad 
\mbox{and } (S_i, E) \reduce \ctrue ,  (S_{i}, P) \Reduceto S_{i + 1} 
\mbox{ for all } 0 \leq i < n 
\\[0.2em] \qquad 
\mbox{and } (S_n, E) \reduce \cfalse 
\end{array}
\]
\end{itemize}
\end{slide}

\end{document}
 

